\linksection{fixt-parametrize}{Fixt. params}

\begin{minted}[autogobble,escapeinside=||,highlightlines=1]{python}
    @pytest.fixture(params=["a.json", "b.json"])
    def config(request: pytest.FixtureRequest):
        return Config(|\highlightbox{request.param}|)

    def test_init(config: Config):
        print(config.filename)
        assert False
\end{minted}
% VS Code workaround
\if 0
    )
    \end{minted}\fi

\vspace{-.5em}
Now each test using the \mono{config} fixture will run twice:
\vspace{.5em}

\begin{centering}
    \begin{tikzpicture}
        \node [fixture, dashed, minimum width=1.5cm] (config) {\mono{config}};
        \node [fixture, below=2mm of config, xshift=-2cm, align=center, minimum width=2cm, minimum height=1cm] (gt) {\mono{config} \\ \py{"a.json"}};
        \node [fixture, below=2mm of config, xshift=2cm, align=center, minimum width=2cm, minimum height=1cm] (rpn) {\mono{config} \\ \py{"b.json"}};

        \node [test, below=4mm of gt] (testgt) {\monot{test\_init[\textcolor{pygnum}{a.json}]}};
        \node [test, below=4mm of rpn] (testrpn) {\monot{test\_init[\textcolor{pygnum}{b.json}]}};

        \draw [-{to}, dashed] (config.west) -| (gt);
        \draw [-{to}, dashed] (config.east) -| (rpn);

        \draw [->] (testgt) -- (gt);
        \draw [->] (testrpn) -- (rpn);
    \end{tikzpicture}
\end{centering}
\vspace{-.5em}

\linksection{indirect}{Indirect parametrization}
\begin{minted}[autogobble,escapeinside=||,highlightlines=7]{python}
    @pytest.fixture
    def config(request):
        return Config(request.param)

    @pytest.mark.parametrize("config, debug", [
        ("prod.json", False), ("dev.json", True),
    ], indirect=["config"])  # or =True for all
    def test_debug(config: Config, debug: bool):
        assert config.debug == debug
\end{minted}

Now the test decides how the fixture gets parametrized,
and the test function gets the fixture return value:
\vspace{.5em}

\begin{tikzpicture}[test/.append style={minimum width=4.5cm}]
    \node [test, dashed, minimum width=5.5cm] (testbase) {\mono{test_debug}};
    \node [test, below=5mm of testbase, minimum height=1cm, minimum width=5.5cm] (testprod) {\monot{test\_debug[\textcolor{pygnum}{prod.json}-\textcolor{pygnum}{False}]}};
    \node [test, below=-0.5pt of testprod, minimum height=1cm, minimum width=5.5cm] (testdev) {\monot{test\_debug[\textcolor{pygnum}{dev.json}-\textcolor{pygnum}{True}]}};

    \node [fixture, right=5mm of testprod, align=center, minimum height=1cm, minimum width=2.3cm] (fixtprod) {\mono{config} \\ \py{"prod.json"}};
    \node [fixture, right=5mm of testdev, align=center, minimum height=1cm, minimum width=2.3cm] (fixtdev) {\mono{config} \\ \py{"dev.json"}};

    \draw [-{to},dashed,gray] (testbase) -- (testprod);
    \draw [->] (testprod.east) -- (fixtprod.west);
    \draw [->] (testdev.east) -- (fixtdev.west);
\end{tikzpicture}

%Can take a list of argument names \\
%(like \py{indirect=["config"]} here), \\
%or \py{indirect=True} to make all indirect.
